\chapter{What is proved}\label{ch:lean}
% src: formal/twin-containment/TwinContainment.lean
% src: formal/twin-containment/README.md
% src: formal/twin-containment/Axioms.lean
% Lean excerpts below are verbatim from TwinContainment.lean. The style maps the file's Unicode
% symbols to math so that pdflatex can typeset them.
\lstdefinestyle{lean}{basicstyle=\ttfamily\scriptsize, columns=fullflexible, breaklines=true,
  frame=single, framerule=0.3pt, extendedchars=true, keepspaces=true,
  literate={ℕ}{{$\mathbb{N}$}}1 {∀}{{$\forall$}}1 {∃}{{$\exists$}}1 {∈}{{$\in$}}1
    {∧}{{$\wedge$}}1 {∨}{{$\vee$}}1 {≤}{{$\le$}}1 {≠}{{$\neq$}}1 {¬}{{$\neg$}}1
    {→}{{$\to$}}1 {←}{{$\leftarrow$}}1 {⟨}{{$\langle$}}1 {⟩}{{$\rangle$}}1}

The previous chapters describe how the twin keeps its language models away from authority.
This chapter states what has been proved about that design in the Lean 4 theorem prover
\cite{demoura2021lean4}, with the Mathlib library \cite{mathlib2020}.
The proofs are about a model of the authority path, not about the Python code.
The model is small enough to read in one sitting.
Every output of a language model enters it as an arbitrary input.
The classifier's events, the chooser's proposals, the analyzer's verdict, DEME's vetoes and
rankings, and the redirect to the monitoring centre are all left unconstrained.
Each theorem therefore holds whatever those models say.
% src: formal/twin-containment/TwinContainment.lean:5-8
% src: formal/twin-containment/TwinContainment.lean:351-356

The model instantiates part (iv) of the No Escape theorem of the \erisml{} foundations paper,
which states that reasoning cannot help an agent escape a structural containment
\cite{bond2026noescape}.
% src: formal/twin-containment/README.md:6-7
% src: C:\source\erisml-lib\docs\papers\foundations\no_escape.tex:172-180
That paper argues the theorem in prose and in general terms.
The Lean file proves a concrete version for one robot, its governor and its scene.
The source is one file of 488 lines, \code{TwinContainment.lean}.
% src: formal/twin-containment/TwinContainment.lean (wc -l = 488)
It is pinned to Lean 4 and Mathlib at version v4.32.2.
% src: formal/twin-containment/lean-toolchain:1
% src: formal/twin-containment/lakefile.toml:4-7

Section~\ref{sec:lean-model} defines the model.
Section~\ref{sec:lean-theorems} states the theorems and says what each means for the robot.
Section~\ref{sec:lean-strata} treats the two theorems about strata, whose hypotheses are the
checks the scene loader makes.
Section~\ref{sec:lean-axioms} reports the axiom audit.
Section~\ref{sec:lean-tests} explains how tests tie the model to the code.
Section~\ref{sec:lean-notcovered} lists what the proofs do not cover.
Appendix~\ref{app:theorems} lists every theorem and lemma in the file.

\section{The model has readings, rulings, histories and a gate}\label{sec:lean-model}

\subsection{Readings and witnesses}

A reading is a record of one sensor's report.
It has a name and four Boolean fields.
The fields say whether the sensor is physical, whether its signature verified, whether the
reading is stale and whether it reports an alert.
% src: formal/twin-containment/TwinContainment.lean:31-36
For a reading $r$ we write $\mathsf{phys}(r)$, $\mathsf{att}(r)$, $\mathsf{stale}(r)$ and
$\mathsf{alert}(r)$ for these fields and $\mathsf{name}(r)$ for its name.
A reading counts when all four conditions favour it.
\begin{equation}
  \mathsf{counts}(r) \;=\; \mathsf{phys}(r) \wedge \mathsf{att}(r) \wedge \neg\,\mathsf{stale}(r)
  \wedge \mathsf{alert}(r).
\end{equation}
The number of witnesses in a list of readings $R$ is the number of distinct names among the
readings that count.
\begin{equation}
  W(R) \;=\; \bigl|\,\{\, \mathsf{name}(r) \mid r \in R,\ \mathsf{counts}(r) \,\}\,\bigr|.
\end{equation}
The Lean definitions are two lines.
% src: formal/twin-containment/TwinContainment.lean:39,42 (verbatim)
\begin{lstlisting}[style=lean]
def Reading.counts (r : Reading) : Bool := r.physical && r.attested && !r.stale && r.alert
def witnesses (rs : List Reading) : ℕ := ((rs.filter Reading.counts).map Reading.name).dedup.length
\end{lstlisting}
In the model a name stands for one independent source.
The code groups sensors by the substrate they share (one hub, one body, one bus), and sensors on
one substrate count as one witness.
% src: twin/brain.py:78-79
% src: tests/test_twin_ladder.py:842-848
The model does not represent substrates.
Its name field plays the role of the substrate.

\subsection{Actions, bars and the governor's ruling}

The model has nine actions.
Eight are governed or protective.
They are calling emergency services ($\mathsf{ems}$), physical help ($\mathsf{physical}$),
the privacy override ($\mathsf{privacy}$), driving off an animal ($\mathsf{driveOff}$),
restraint ($\mathsf{restrain}$), the stun device ($\mathsf{stun}$), the spray
($\mathsf{spray}$) and interposing ($\mathsf{interpose}$).
The ninth, $\mathsf{benign}$, stands for every action the scene does not govern.
% src: formal/twin-containment/TwinContainment.lean:62-64
Each action $a$ has a witness bar $\beta(a)$.
\begin{equation}
  \beta(a) \;=\;
  \begin{cases}
    3 & \text{if } a \in \{\mathsf{restrain}, \mathsf{stun}, \mathsf{spray}\},\\
    2 & \text{otherwise}.
  \end{cases}
\end{equation}
% src: formal/twin-containment/TwinContainment.lean:66-69
These values come from the \code{witness\_bar} field of the scene's capability list.
% src: twin/scene/margaret_home.erisml:1168-1170

A ruling is an outcome together with a restraint flag.
The four outcomes are $\mathsf{elevate}$, $\mathsf{authorizeEms}$, $\mathsf{refuseReview}$ and
$\mathsf{refuse}$.
% src: formal/twin-containment/TwinContainment.lean:73-79
The governor's ruling on a request for action $a$ takes the readings $R$ and three flags.
The flag $f$ is the attestation gate's judgement that the signals are fresh.
The flag $c$ is the system's record that the monitoring centre could not be reached.
The flag $k$ is the analyzer's verdict, which comes from a model and is arbitrary.
Write $W = W(R)$.
The outcome of $\mathrm{rule}(a, R, f, c, k)$ is the first case that applies.
\begin{equation}\label{eq:rule}
  \mathrm{outcome} \;=\;
  \begin{cases}
    \mathsf{refuse} & \text{if } \neg f,\\
    \mathsf{elevate} & \text{if } \beta(a) \le W \text{ and } k,\\
    \mathsf{elevate} & \text{if } 2 \le W \text{ and } k,\\
    \mathsf{authorizeEms} & \text{if } a = \mathsf{ems},\ c,\ 1 \le W \text{ and } k,\\
    \mathsf{refuseReview} & \text{if } 1 \le W,\\
    \mathsf{refuse} & \text{otherwise}.
  \end{cases}
\end{equation}
The restraint flag is set only in the second case, and only when $\beta(a) \ge 3$.
% src: formal/twin-containment/TwinContainment.lean:84-90
A request for restraint with two witnesses is therefore elevated without restraint.
This mirrors the governor of the code, \code{govern\_live} in \code{twin/brain.py}.
% src: formal/twin-containment/TwinContainment.lean:81-83
% src: twin/brain.py:339

\subsection{Histories of system events and model events}

The robot's moral state steps on a history, a finite list of events.
The model has seven kinds of event.
\begin{itemize}
  \item A governor ruling, stored with the readings it ruled on.
  \item The owner's opt-in, which arrives on an authenticated channel.
  \item A reflex's force measurement, severe or not.
  \item The world's measure of whether Margaret is inside the spray's reach.
  \item The record that the robot tried the centre and could not reach it.
  \item An action the robot performed.
  \item A model event, which carries an arbitrary string.
\end{itemize}
% src: formal/twin-containment/TwinContainment.lean:117-126
The first six are system events.
The seventh stands for anything a model classified.
Nothing in the model may depend on its content.

The strip of a history removes its model events.
\begin{equation}
  \mathrm{strip}(h) \;=\; [\, e \in h \mid e \text{ is not a model event} \,].
\end{equation}
% src: formal/twin-containment/TwinContainment.lean:166

\subsection{Permitted actions}

The scene's permission for each action is a Boolean function $P(h, a)$ of the history.
Let $\ell(h)$ be the latest ruling in $h$, if there is one.
Calling emergency services is permitted when $\ell(h)$ is an elevation or an authorization of
emergency services.
Physical help, the privacy override and driving off are permitted when $\ell(h)$ is an
elevation.
Restraint is permitted when some ruling in $h$ carried the restraint flag.
The stun device is permitted when $h$ contains the opt-in, a restraint ruling, a severe force
measurement and an interposition.
The spray needs all of these, and in addition the world's latest measure must not place Margaret
inside its reach.
Interposing and the benign actions are always permitted.
% src: formal/twin-containment/TwinContainment.lean:154-162
The Lean definition reads as follows.
\begin{lstlisting}[style=lean]
def permitted (h : List Ev) : Action → Bool
  | .ems => latestIs .elevate h || latestIs .authorizeEms h
  | .physical | .privacy | .driveOff => latestIs .elevate h
  | .restrain => has isRestraint h
  | .stun => has isOptIn h && has isRestraint h && has isSevere h && has isInterposed h
  | .spray => has isOptIn h && has isRestraint h && has isSevere h && has isInterposed h && !inZone h
  | .interpose | .benign => true
\end{lstlisting}
% src: formal/twin-containment/TwinContainment.lean:156-162 (verbatim)
With no zone measure in the history the model treats Margaret as outside the spray's reach.
% src: formal/twin-containment/TwinContainment.lean:148
The scene has the same default.
Its spray prohibition \code{ll5spray} fires only when the latest zone event says Margaret is
inside.
% src: twin/scene/margaret_home.erisml:717-725
% src: twin/scene/margaret_home.erisml:797

\subsection{Reachable histories}

The model says how the twin runs through an inductive predicate $\mathrm{Reach}$ on histories.
We write $h \cdot e$ for the history $h$ followed by the event $e$.
The empty history is reachable.
If $h$ is reachable, so is each of the following extensions.
\begin{enumerate}
  \item $h \cdot \mathsf{model}(s)$ for any string $s$.
  Models may add any event at any time.
  \item $h \cdot \mathsf{ruling}(\mathrm{rule}(a, R, f, c, k), R)$ for any action, readings and
  flags, provided that $c$ is true only when $h$ already records the centre as unreachable.
  A ruling is only ever the governor's output.
  \item $h \cdot \mathsf{optIn}$, $h \cdot \mathsf{attackMeasured}(b)$,
  $h \cdot \mathsf{sprayZone}(b)$ and $h \cdot \mathsf{centreUnavailable}$.
  \item $h \cdot \mathsf{performed}(a)$, provided that $P(h, a)$ holds.
  The robot executes an action only when the scene permits it.
\end{enumerate}
% src: formal/twin-containment/TwinContainment.lean:200-212
The readings $R$ in a ruling are arbitrary.
The model does not assume they report the world truly.
It assumes only that the attested flag of each reading is correct, which the code's signature
check establishes (Sections~\ref{sec:cont-attest} and~\ref{sec:lean-tests}).
Figure~\ref{fig:lean-model} shows the objects of the model and the path a run takes through
them.

\begin{figure}[t]
  \centering
  % Chapter 10. The objects of the Lean model and how a run moves through them.
% src: formal/twin-containment/TwinContainment.lean:31-42 (readings, witnesses), 66-90 (bar, rule),
%      117-166 (events, permitted, strip), 203-212 (Reach), 357-366 (gate)
\begin{tikzpicture}[house,
  p/.style={minimum width=23mm, minimum height=46mm, text width=21mm, anchor=north},
  d/.style={detail, text width=19mm}]
\foreach \i/\col/\x in {1/Purple/0, 2/Blue/28, 3/Teal/56, 4/Green/84, 5/Amber/112, 6/Orange/140}{
  \node[panel=\col, p] (P\i) at (\x mm,0) {};
  \node[step=\col] at ([yshift=-5mm]P\i.north) {\i};
}
\node[stepname=Purple, below=10mm of P1.north] (n1) {Readings};
\node[desc, below=0mm of n1, text width=21mm] (e1) {one per sensor};
\node[d, below=1.5mm of e1] {name\\physical\\attested\\stale\\alert};

\node[stepname=Blue, below=10mm of P2.north] (n2) {Witnesses};
\node[desc, below=0mm of n2, text width=21mm] (e2) {$W(R)$};
\node[d, below=1.5mm of e2] {count a reading\\if physical\\attested, fresh\\and alerting\\distinct names};

\node[stepname=Teal, below=10mm of P3.north] (n3) {Ruling};
\node[desc, below=0mm of n3, text width=21mm] (e3) {$\mathrm{rule}(a,R,f,c,k)$};
\node[d, below=1.5mm of e3] {bar 2, or 3 for\\restraint and\\devices\\1 for EMS when\\centre is down};

\node[stepname=Green, below=10mm of P4.north] (n4) {History};
\node[desc, below=0mm of n4, text width=21mm] (e4) {list of events};
\node[d, below=1.5mm of e4] {rulings, opt-in\\force, zone\\centre down\\performed\\model events};

\node[stepname=Amber, below=10mm of P5.north] (n5) {Permitted};
\node[desc, below=0mm of n5, text width=21mm] (e5) {$P(h,a)$};
\node[d, below=1.5mm of e5] {reads system\\events only\\so equal on $h$\\and on\\$\mathrm{strip}(h)$};

\node[stepname=Orange, below=10mm of P6.north] (n6) {Gate};
\node[desc, below=0mm of n6, text width=21mm] (e6) {one action out};
\node[d, below=1.5mm of e6] {redirect, proposal\\first ranked\\or benign\\always permitted};

% the authority path
\draw[flow] ([yshift=-23mm]P1.north east) -- ([yshift=-23mm]P2.north west);
\draw[grant] ([yshift=-23mm]P2.north east) -- ([yshift=-23mm]P3.north west);
\draw[grant] ([yshift=-23mm]P3.north east) -- ([yshift=-23mm]P4.north west);
\draw[flow] ([yshift=-23mm]P4.north east) -- ([yshift=-23mm]P5.north west);
\draw[flow] ([yshift=-23mm]P5.north east) -- ([yshift=-23mm]P6.north west);
% the performed action returns to the history
\draw[flow] (P6.south) -- ++(0,-5mm) -| node[flowlabel, pos=0.25, below]{performed action appended to the history} (P4.south);

% the models, whose outputs are arbitrary inputs
\node[panel=Grey, text width=95mm, minimum height=9mm, anchor=south] (M) at ([yshift=9mm]$(P3.north)!0.5!(P6.north)$)
  {\textbf{Language models}\\analyzer verdict, classified events, DEME vetoes and ranking, redirect, proposal};
\draw[untrusted] (P3.north |- M.south) -- node[flowlabel, right]{$k$} (P3.north);
\draw[untrusted] (P4.north |- M.south) -- node[flowlabel, right]{model events} (P4.north);
\draw[untrusted] (P6.north |- M.south) -- node[flowlabel, right]{$p, V, L, d$} (P6.north);
\end{tikzpicture}

  \caption{The objects of the Lean model in the order a run uses them. Readings become witnesses,
  witnesses meet a bar in the governor's ruling, rulings enter the history, the history fixes what
  is permitted and the gate picks one permitted action. Red edges carry authority and dashed
  edges carry model outputs. The analyzer's verdict can only block a ruling, and the permission
  function ignores every model event in the history.}
  \label{fig:lean-model}
\end{figure}

\subsection{The DEME output gate}

The output gate of Section~\ref{sec:cont-gate} sits between the brain's proposal and the body.
In the model the gate takes the history $h$, the proposal $p$, DEME's veto function $V$,
DEME's ranked list $L$ and an optional redirect $d$.
The redirect is the monitoring centre, chosen when the tragic-conflict module flags a deliberate
decision.
% src: formal/twin-containment/TwinContainment.lean:351-356
All of $V$, $L$ and $d$ are arbitrary.
The gate returns the first of the following that the scene permits and DEME does not veto.
\begin{enumerate}
  \item The redirect $d$, if there is one.
  \item The proposal $p$.
  \item The first action in $L$ that the scene permits and DEME does not veto.
  \item Failing all of these, the benign idle action, which the scene always permits.
\end{enumerate}
% src: formal/twin-containment/TwinContainment.lean:357-366

\subsection{Strata}

A stratification has a state space $S$ and an event space $E$.
A semantic gate $g$ has a source predicate $\mathrm{src}_g$ on $S$, a firing predicate
$\mathrm{fires}_g$ on $E$ and a destination $\mathrm{dst}_g \in S$.
% src: formal/twin-containment/TwinContainment.lean:434-437
Given a list of gates $G$, one step on event $e$ from state $s$ moves to the destination of the
first gate in $G$ whose source admits $s$ and which fires on $e$.
If no gate matches, the state stays where it is.
\begin{equation}
  \mathrm{sstep}_G(s, e) \;=\;
  \begin{cases}
    \mathrm{dst}_g & \text{if } g \text{ is the first gate in } G \text{ with }
      \mathrm{src}_g(s) \wedge \mathrm{fires}_g(e),\\
    s & \text{if there is no such gate}.
  \end{cases}
\end{equation}
A run folds the step over a list of events,
$\mathrm{srun}_G(s, [e_1, \dots, e_n]) = \mathrm{sstep}_G(\cdots \mathrm{sstep}_G(s, e_1) \cdots, e_n)$.
% src: formal/twin-containment/TwinContainment.lean:439-445
The runtime does the same.
Its \code{cross} method takes the first gate that fires and moves to its target.
% src: C:\source\erisml-compiler-monitor\src\erisml_compiler\runtime\strata.py:124-148
The state space and the gates are left general, so these definitions cover every stratification
a scene can declare.

\section{Each granting outcome rests on attested witnesses}\label{sec:lean-theorems}

The file proves 28 theorems.
The file \code{Axioms.lean} audits 22 of them.
Five of the other six are helper lemmas used inside other proofs.
The sixth is \code{unattested\_never\_counts}.
% src: formal/twin-containment/TwinContainment.lean (theorem declarations counted)
% src: formal/twin-containment/Axioms.lean:3-24 (the 22 audited results)
This section groups the results by what they protect.
Throughout, $a$ is any action, $R$ any readings and $f, c, k$ any flags.

\subsection{Junk readings add no witness}

\begin{description}
  \item[\code{junk\_adds\_no\_witness}.] If no reading in a list $J$ counts, then
  $W(R \mathbin{+\!\!+} J) = W(R)$.
  \item[\code{unattested\_never\_counts}.] If $\neg\,\mathsf{att}(r)$ then
  $\neg\,\mathsf{counts}(r)$.
  \item[\code{repeat\_adds\_no\_witness}.] $W(r :: r :: R) = W(r :: R)$.
\end{description}
% src: formal/twin-containment/TwinContainment.lean:44-58
For the robot these say that a forged reading, a stale one, a network message and a repeated
reading add nothing to the count.
An attacker who can inject messages but cannot sign as a physical sensor cannot raise the
witness count by volume.

\subsection{The analyzer can only refuse}

\begin{description}
  \item[\code{analyzer\_can\_only\_refuse}.] With $k$ false, the outcome of
  $\mathrm{rule}(a, R, f, c, k)$ is neither $\mathsf{elevate}$ nor $\mathsf{authorizeEms}$.
\end{description}
% src: formal/twin-containment/TwinContainment.lean:92-95
The analyzer is a model.
Its verdict appears in the ruling only as a conjunct of the granting cases of
equation~\eqref{eq:rule}.
A negative verdict therefore blocks a grant, and a positive verdict alone grants nothing.

\subsection{Each grant meets its bar}

\begin{description}
  \item[\code{elevate\_needs\_two}.] If the outcome of $\mathrm{rule}(a, R, f, c, k)$ is
  $\mathsf{elevate}$ then $W(R) \ge 2$.
  \item[\code{restraint\_needs\_three}.] If the ruling carries the restraint flag then
  $W(R) \ge 3$ and the outcome is $\mathsf{elevate}$.
  \item[\code{authorize\_ems\_only\_when\_ladder\_exhausted}.] If the outcome is
  $\mathsf{authorizeEms}$ then $a = \mathsf{ems}$, $c$ holds and $W(R) \ge 1$.
  \item[\code{nothing\_without\_a\_witness}.] If $W(R) = 0$ the outcome is neither
  $\mathsf{elevate}$ nor $\mathsf{authorizeEms}$.
\end{description}
% src: formal/twin-containment/TwinContainment.lean:97-113
These are facts about one ruling.
Two independent attested physical sensors are needed to elevate.
Three are needed to restrain or to use a device.
One is enough only for a call to emergency services, and only when the system has recorded that
the centre could not be reached.
The proofs unfold the definition of the ruling and close each case by arithmetic.
% src: formal/twin-containment/TwinContainment.lean:99-100

\subsection{Model events cannot change a governed permission}

\begin{description}
  \item[\code{no\_model\_influence}.] For every history $h$ and action $a$,
  $P(\mathrm{strip}(h), a) = P(h, a)$.
  \item[\code{reasoning\_cannot\_help}.] If $\mathrm{strip}(h) = \mathrm{strip}(h')$ then
  $P(h, a) = P(h', a)$ for every $a$.
\end{description}
% src: formal/twin-containment/TwinContainment.lean:184-196
This is a non-interference property in the sense of information flow.
Two histories that agree on their system events permit the same actions, whatever the models
classified in either.
The Lean statement of the first is one line.
% src: formal/twin-containment/TwinContainment.lean:184
\begin{lstlisting}[style=lean]
theorem no_model_influence (h : List Ev) (a : Action) : permitted (strip h) a = permitted h a := by
\end{lstlisting}
Its proof rests on two lemmas.
\code{filterMap\_strip} says that any extraction which ignores model events gives the same
result on $h$ and on $\mathrm{strip}(h)$.
\code{any\_strip} says the same for any test of whether some event has a property.
% src: formal/twin-containment/TwinContainment.lean:168-182
The permission function reads the history only through such extractions and tests.
For the robot this means that no sentence a model produces, no event it invents and no event it
suppresses changes what the scene permits.
A model can still choose among permitted actions.
It cannot widen the set.

\subsection{Every executed action rests on its evidence}

The theorems so far concern a single ruling or a single history.
The next five concern every reachable history.
\begin{description}
  \item[\code{rulings\_are\_governed}.] If $\mathrm{Reach}(h)$ and
  $\mathsf{ruling}(r, R) \in h$, then $r = \mathrm{rule}(a, R, f, c, k)$ for some
  $a, f, c, k$.
  \item[\code{ems\_rests\_on\_evidence}.] If $\mathrm{Reach}(h)$ and $P(h, \mathsf{ems})$, then
  the latest ruling $r$ of $h$ is recorded in $h$ with readings $R$ and $W(R) \ge 1$.
  \item[\code{elevated\_rests\_on\_two}.] If $\mathrm{Reach}(h)$ and $P(h, a)$ for
  $a \in \{\mathsf{physical}, \mathsf{privacy}, \mathsf{driveOff}\}$, then $h$ contains a ruling
  with outcome $\mathsf{elevate}$ on readings $R$ with $W(R) \ge 2$.
  \item[\code{restraint\_rests\_on\_three}.] If $\mathrm{Reach}(h)$ and $P(h, a)$ for
  $a \in \{\mathsf{restrain}, \mathsf{stun}, \mathsf{spray}\}$, then $h$ contains a ruling with the
  restraint flag on readings $R$ with $W(R) \ge 3$.
  \item[\code{executed\_were\_permitted}.] If $\mathrm{Reach}(h)$ and
  $\mathit{pre} \cdot \mathsf{performed}(a)$ is a prefix of $h$, then $P(\mathit{pre}, a)$.
\end{description}
% src: formal/twin-containment/TwinContainment.lean:214-347
The first is proved by induction on $\mathrm{Reach}$.
Only the ruling step adds a ruling, and it adds the governor's.
The next three combine it with the per-ruling bars.
The last is also by induction, through the lemma \code{prefix\_snoc}, which splits a prefix of
$h \cdot x$ into a prefix of $h$ or the whole of it.
% src: formal/twin-containment/TwinContainment.lean:306-312

Together they give the main guarantee for the robot.
In every run the model allows, each governed action the robot executed was permitted on the
history before it.
That permission came from a governor ruling resting on at least as many attested physical
witnesses as the action's bar.
The bar is one for emergency services when the centre was unreachable, two for the other
elevated actions and three for restraint and the two devices.
% src: formal/twin-containment/TwinContainment.lean:14-17
The statement for restraint reads as follows in Lean.
% src: formal/twin-containment/TwinContainment.lean:282-283
\begin{lstlisting}[style=lean]
theorem restraint_rests_on_three (h : List Ev) (hr : Reach h) (a : Action) (ha : a = .restrain ∨ a = .stun ∨ a = .spray)
    (hp : permitted h a = true) : ∃ rs r, Ev.ruling r rs ∈ h ∧ r.restraint = true ∧ 3 ≤ witnesses rs := by
\end{lstlisting}
Figure~\ref{fig:lean-proofs} shows which lemmas each of these results uses.

\begin{figure}[t]
  \centering
  % Chapter 10. Which lemmas each result of TwinContainment.lean uses in its proof.
% src: formal/twin-containment/TwinContainment.lean (proof bodies: 99, 112, 185-190, 196, 271-278,
%      290-291, 301-304, 321-347, 387, 393, 404)
\begin{tikzpicture}[house,
  lem/.style={detail, font=\ttfamily\scriptsize, text=hsText, text width=34mm, minimum height=6mm},
  res/.style={lem, draw=hsGreyLine, line width=0.5pt},
  dep/.style={flow, line width=1.1pt}]
% one ruling
\node[lem] (tlb) at (12mm, 14mm) {two\_le\_bar};
\node[res] (nww) at (12mm, 26mm) {nothing\_without\_\\a\_witness};
\node[res] (ent) at (24mm, 0mm) {elevate\_needs\_two};
\node[res] (aem) at (24mm, -13mm) {authorize\_ems\_only\_\\when\_ladder\_exhausted};
\node[res] (rnt) at (24mm, -26mm) {restraint\_needs\_three};
% every reachable run
\node[res] (ers) at (80mm, 0mm) {elevated\_rests\_on\_two};
\node[res] (eme) at (80mm, -13mm) {ems\_rests\_on\_evidence};
\node[res] (rrt) at (80mm, -26mm) {restraint\_rests\_on\_three};
\node[res] (ewp) at (80mm, -39mm) {executed\_were\_permitted};
\node[res] (rag) at (130mm, -13mm) {rulings\_are\_governed};
\node[lem] (lm) at (130mm, 4mm) {latest\_mem};
\node[lem] (psn) at (130mm, -39mm) {prefix\_snoc};

\begin{scope}[on background layer]
  \node[panel=Blue, fit=(tlb)(nww)(ent)(aem)(rnt), inner sep=3mm, label={[stepname=Blue]north:One ruling}] {};
  \node[panel=Purple, fit=(ers)(ewp)(rag)(lm)(psn), inner sep=3mm, label={[stepname=Purple]north:Every reachable run}] {};
\end{scope}

\draw[dep] (tlb) -- (nww);
\draw[dep] (tlb) -- (ent);
\draw[dep] (ent) -- (ers);
\draw[dep] (ent) -- (eme);
\draw[dep] (aem) -- (eme);
\draw[dep] (rnt) -- (rrt);
\draw[dep] (rag) -- (ers);
\draw[dep] (rag) -- (eme);
\draw[dep] (rag) -- (rrt);
\draw[dep] (lm) -- (ers);
\draw[dep] (lm) -- (eme);
\draw[dep] (psn) -- (ewp);

% non-interference and the gate
\node[lem] (fms) at (12mm, -66mm) {filterMap\_strip};
\node[lem] (ans) at (12mm, -78mm) {any\_strip};
\node[res] (nmi) at (50mm, -72mm) {no\_model\_influence};
\node[res] (rch) at (50mm, -86mm) {reasoning\_cannot\_help};
\node[res] (fbp) at (100mm, -66mm) {fallback\_permitted};
\node[res] (gpm) at (138mm, -66mm) {gate\_permitted};
\node[res] (gsr) at (138mm, -80mm) {gated\_step\_reachable};
\node[res] (gps) at (100mm, -80mm) {gate\_passes};
\begin{scope}[on background layer]
  \node[panel=Teal, fit=(fms)(ans)(nmi)(rch), inner sep=3mm, label={[stepname=Teal]north:Non-interference}] {};
  \node[panel=Green, fit=(fbp)(gpm)(gsr)(gps), inner sep=3mm, label={[stepname=Green]north:Output gate}] {};
\end{scope}
\draw[dep] (fms) -- (nmi);
\draw[dep] (ans) -- (nmi);
\draw[dep] (nmi) -- (rch);
\draw[dep] (fbp) -- (gpm);
\draw[dep] (gpm) -- (gsr);

% results proved directly from the definitions
\node[panel=Grey, text width=150mm, anchor=north, label={[stepname=Grey]north:Proved from the definitions alone}] at (75mm, -100mm)
  {\ttfamily\scriptsize junk\_adds\_no\_witness\quad unattested\_never\_counts\quad repeat\_adds\_no\_witness\quad analyzer\_can\_only\_refuse\\[2pt]
   dispatch\_ignores\_the\_robot\quad forged\_telemetry\_is\_no\_hazard\quad authority\_needs\_system\quad absorbing\_stays};
\end{tikzpicture}

  \caption{The lemmas each result of the Lean file uses in its proof. Outlined boxes are
  results and plain boxes are helper lemmas. The guarantees about every reachable run are built
  from the bars of a single ruling and from \code{rulings\_are\_governed}, which says that every
  ruling in a run is the governor's.}
  \label{fig:lean-proofs}
\end{figure}

\subsection{The output gate keeps the robot inside the permitted set}

\begin{description}
  \item[\code{gate\_permitted}.] For every history, proposal, veto function, ranking and
  redirect, the output $o = \mathrm{gate}(h, p, V, L, d)$ satisfies $P(h, o)$.
  \item[\code{fallback\_permitted}.] The fallback choice, the first ranked action the scene
  permits and DEME does not veto or else the benign action, is permitted.
  \item[\code{gate\_passes}.] Without a redirect, if $P(h, p)$ and $V(p)$ is false, the gate
  returns $p$.
  \item[\code{gated\_step\_reachable}.] If $\mathrm{Reach}(h)$ then
  $\mathrm{Reach}(h \cdot \mathsf{performed}(\mathrm{gate}(h, p, V, L, d)))$.
\end{description}
% src: formal/twin-containment/TwinContainment.lean:368-404
The first says that DEME cannot make the robot leave the permitted set.
Whatever it vetoes, however it ranks and wherever the redirect points, the gate's output is
permitted by the scene.
The third says that DEME can only veto.
It cannot replace a permitted, unvetoed proposal when there is no redirect.
The fourth joins the gate to the run model.
A step through the gate is a reachable step, so every theorem about reachable histories covers
runs that go through the gate.

\subsection{The centre's dispatch ignores the robot}

The monitoring centre is a second agent with its own scene.
Its dispatch of emergency services is permitted when the telemetry contains a reading that
counts, when Margaret did not answer on the speaker line, or when her own words, read alone,
ask for help.
% src: formal/twin-containment/TwinContainment.lean:408-414
The robot's message is an argument of the permission function and is unused.
\begin{description}
  \item[\code{dispatch\_ignores\_the\_robot}.] For any telemetry and flags, the dispatch
  permission is the same for every two robot messages $m$ and $m'$.
  \item[\code{forged\_telemetry\_is\_no\_hazard}.] If no reading in the telemetry is attested,
  the telemetry shows no hazard.
\end{description}
% src: formal/twin-containment/TwinContainment.lean:416-422
The first is proved by reflexivity.
% src: formal/twin-containment/TwinContainment.lean:416-417
\begin{lstlisting}[style=lean]
theorem dispatch_ignores_the_robot (tele : List Reading) (n w : Bool) (m m' : String) :
    dispatchPermitted tele n w m = dispatchPermitted tele n w m' := rfl
\end{lstlisting}
The proof is trivial because the property is built into the definition.
What matters is that the definition matches the centre's scene, which the static tests check
(Section~\ref{sec:lean-tests}).
A robot that is fooled, or a model that writes an alarming message, cannot by that message
alone send an ambulance.

\section{The strata theorems assume exactly the loader's checks}\label{sec:lean-strata}

Section~\ref{sec:cont-strata} describes the seven strata of the home.
Two theorems cover every stratification at once.
Let $\mathrm{auth}$ be a predicate on states marking the authority strata, those that grant
something.
Let $\mathrm{sys}$ be a predicate on events marking system events.
\begin{description}
  \item[\code{authority\_needs\_system}.] Suppose every gate whose destination is an authority
  state fires only on system events.
  Then for every list of events none of which is a system event, and every start state $s$ with
  $\neg\,\mathrm{auth}(s)$, the final state $\mathrm{srun}_G(s, es)$ is not an authority state.
\end{description}
% src: formal/twin-containment/TwinContainment.lean:447-465
Let $\mathrm{absorb}$ mark the absorbing strata and $\mathrm{ovs}$ the oversight events.
\begin{description}
  \item[\code{absorbing\_stays}.] Suppose every gate that can fire from an absorbing state to a
  non-absorbing one fires only on oversight events.
  Then for every list of events none of which is an oversight event, and every absorbing start
  state $s$, the final state $\mathrm{srun}_G(s, es)$ is absorbing.
\end{description}
% src: formal/twin-containment/TwinContainment.lean:467-486
Both proofs are by induction on the event list, generalizing the start state.
% src: formal/twin-containment/TwinContainment.lean:453,474
Absorption follows Definition 8.6 of Geometric Ethics, chapter 8~\cite{bond2026geometricethics}.
% src: formal/twin-containment/TwinContainment.lean:467

The hypotheses on the gates are what the \erisml{} compiler's scene loader checks before a scene
may run.
% src: formal/twin-containment/TwinContainment.lean:426-432
The loader's method \code{Stratification.\_check} raises an error for any gate into an authority
state whose triggering event type is not declared with source \code{system}.
It also raises an error for any gate whose source states include an absorbing state other than
its target, unless its trigger is a declared oversight event.
% src: C:\source\erisml-compiler-monitor\src\erisml_compiler\runtime\strata.py:105-121
The second check is slightly stronger than the Lean hypothesis.
It also refuses a gate from one absorbing state to another without oversight.
It therefore establishes the hypothesis.
A third check refuses a stratification whose initial state is an authority state.
% src: C:\source\erisml-compiler-monitor\src\erisml_compiler\runtime\strata.py:84-85
That is the hypothesis $\neg\,\mathrm{auth}(s)$ for the initial state.
A scene that breaks any of these rules never runs.
% src: C:\source\erisml-compiler-monitor\src\erisml_compiler\runtime\strata.py:26-33
Figure~\ref{fig:lean-strata} sets the hypotheses beside the checks.

\begin{figure}[t]
  \centering
  % Chapter 10. The two strata theorems: their gate hypotheses and the loader checks that establish them.
% src: formal/twin-containment/TwinContainment.lean:447-486 (authority_needs_system, absorbing_stays)
% src: C:\source\erisml-compiler-monitor\src\erisml_compiler\runtime\strata.py:84-85,105-121
\begin{tikzpicture}[house,
  d/.style={detail, text width=60mm, align=left}]
% a generic stratification
\node[stratum] (s0) at (0,0) {initial};
\node[stratum] (s1) at (34mm,0) {other};
\node[authority] (au) at (0,-24mm) {authority};
\node[absorbing] (ab) at (34mm,-24mm) {absorbing};
\draw[untrusted] (s0) -- node[flowlabel, above]{any event} (s1);
\draw[grant] (s0) -- node[flowlabel, left, align=right]{system\\event only} (au);
\draw[flow] ([xshift=-3mm]s1.south) -- node[flowlabel, left, align=right]{any\\event} ([xshift=-3mm]ab.north);
\draw[flow] ([xshift=3mm]ab.north) -- node[flowlabel, right, align=left]{oversight\\event only} ([xshift=3mm]s1.south);
\draw[untrusted] (s1.south west) -- node[pos=0.5, text=hsRed, font=\sffamily\bfseries\normalsize]{$\times$} (au.north east);
\node[figsub, align=center] at (17mm,-34mm) {a model event into an authority state\\is refused when the scene loads};

% the correspondence
\node[stepname=Teal, anchor=north west] (Lt) at (64mm, 6mm) {Lean hypothesis};
\node[d, below=1.5mm of Lt.south west, anchor=north west] (L1) {\textbf{authority\_needs\_system}\\every gate into an authority state fires only on a system event, and the start state is not an authority state};
\node[d, below=1.5mm of L1] (L2) {\textbf{absorbing\_stays}\\every gate that can leave an absorbing state fires only on an oversight event};
\node[stepname=Blue, below=12mm of L2.south west, anchor=north west] (Ct) {Loader check before a scene runs};
\node[d, below=1.5mm of Ct.south west, anchor=north west] (C1) {refuses a gate into an authority state whose trigger is not declared with source system, and an authority state as the initial state};
\node[d, below=1.5mm of C1] (C2) {refuses a gate that leaves an absorbing state on an event not declared as oversight};
\begin{scope}[on background layer]
  \node[panel=Teal, fit=(Lt)(L1)(L2), inner sep=2.5mm] (L) {};
  \node[panel=Blue, fit=(Ct)(C1)(C2), inner sep=2.5mm] (C) {};
\end{scope}
\draw[flow] (C.north) -- node[flowlabel, right]{establishes} (L.south);
\end{tikzpicture}

  \caption{The two strata theorems on a generic stratification, with their hypotheses beside
  the loader checks that establish them. The only edge into the authority state carries a system
  event, and the only edge out of the absorbing state carries an oversight event.}
  \label{fig:lean-strata}
\end{figure}

For the robot this means the two theorems hold for every scene that loads.
The twin's \code{visitor\_standing} stratification is one such scene element.
Its authority states are \code{arranged} and \code{household}, and its absorbing state is
\code{hostile}.
% src: twin/scene/margaret_home.erisml:816-820
No sequence of model-classified events can make a visitor arranged or household.
A visitor's claim about himself moves nothing.
A visitor who attacked stays hostile until an oversight event, the centre's \code{visitor\_cleared},
moves him out.
% src: twin/scene/margaret_home.erisml:802-804
% src: tests/test_twin_ladder.py:1497 (test_what_a_visitor_says_about_himself_moves_nothing)
% src: tests/test_twin_ladder.py:1508 (test_a_visitor_who_attacked_stays_hostile_until_the_centre_clears_him)
The theorems say more than the tests do.
The tests try particular traces.
The theorems cover every trace.

\section{The proofs use only Lean's standard axioms}\label{sec:lean-axioms}

The file \code{Axioms.lean} asks Lean to print the axioms on which each of the 22 results
depends.
% src: formal/twin-containment/Axioms.lean:1-24
The reported set is Lean's three standard axioms.
They are propositional extensionality (\code{propext}), the axiom of choice
(\code{Classical.choice}) and the soundness of quotients (\code{Quot.sound}).
% src: formal/twin-containment/README.md:28-29
No theorem uses \code{sorry}, the placeholder that admits an unproved statement.
% src: formal/twin-containment/README.md:28 ("All 22 are proved without sorry")
% src: TwinContainment.lean contains no occurrence of "sorry" (grep, 2026-10-03)
No axiom is declared in the file.
The commit records state that the file was built on Atlas with these results.
% src: git log of formal/twin-containment, commits 2b39d1d and 7f2625a
\TODO{archive the lake build log and the Axioms.lean output with the monograph's artifacts}

Lean's report for a theorem includes the axioms of every lemma that theorem uses.
The five helper lemmas are used by audited results, so their axioms are covered.
The remaining theorem, \code{unattested\_never\_counts}, is neither printed nor used by a
printed result.
% src: formal/twin-containment/Axioms.lean (no line for unattested_never_counts)
Its proof is a single simplification of the definition of $\mathsf{counts}$.
% src: formal/twin-containment/TwinContainment.lean:50-51
\TODO{add unattested\_never\_counts to Axioms.lean and record its output}

The build is reproducible from the pinned toolchain.
One runs \code{lake build} for the proofs and \code{lake env lean Axioms.lean} for the audit.
% src: formal/twin-containment/README.md:33

\section{Static tests tie the model to the code}\label{sec:lean-tests}

The proofs are about the model.
Whether the Python code and the scene files implement the model is a separate question.
Two kinds of test address it.
% src: formal/twin-containment/README.md:22-26

The first kind checks the brain and the compiled scene at each boundary the model draws.
The governor test runs \code{govern\_live} on two attested witnesses, one, none, forged devices
and a reachable or unreachable centre.
It checks the outcome and the bar against the cases of equation~\eqref{eq:rule}.
% src: tests/test_twin_ladder.py:309-333 (test_the_governor_bar)
Other tests check that a signature by the wrong key, a payload altered after signing, a replayed
reading and a stale reading do not count.
% src: tests/test_twin_ladder.py:803-840
These are what justify the model's assumption that the attested flag is correct.
A further test checks that a fooled operator model at the centre cannot dispatch on the robot's
claim, and that unsigned telemetry is no ground.
% src: tests/test_twin_ladder.py:638-662
% src: tests/test_twin_ladder.py:681-689

The second kind is structural.
These tests read the real scene files and require them to have the shape the model assumes.
The test file names them as the bridge to the proof.
% src: tests/test_twin_ladder.py:705-709
If someone adds a model-classified condition to a governed prohibition, the continuous
integration run fails.
There are five structural tests.
\begin{enumerate}
  \item Every governed robot action is denied by default.
  Each is guarded by a prohibition, directly or through the general guard on elevated actions.
  % src: tests/test_twin_ladder.py:749-754
  \item Every prohibition guarding a governed action rests only on system events.
  Duty-of-care prohibitions are exempt because they only add restrictions.
  % src: tests/test_twin_ladder.py:757-768
  \item No permission norm grants a governed action.
  % src: tests/test_twin_ladder.py:771-774
  \item The centre's dispatch guard rests only on system events.
  % src: tests/test_twin_ladder.py:777-785
  \item Every condition that lifts or restores privacy rests on system events, and every
  oversight event is a system event.
  % src: tests/test_twin_ladder.py:788-795
\end{enumerate}
A token rests on system events when it names an event of a type declared with source
\code{system}, a governor ruling or a performed action.
A condition token qualifies when its own tokens do, and a stratum token qualifies when the
compiler's token checker accepts it.
% src: tests/test_twin_ladder.py:720-742 (_tokens_are_structural)
The first three tests are what make \code{permitted} a faithful summary of the scene.
The fourth is what makes \code{dispatchPermitted} faithful to the centre's scene.
Figure~\ref{fig:lean-bridge} lists each tie.

\begin{figure}[t]
  \centering
  % Chapter 10. How each part of the Lean model is tied to the code and the scene files.
% src: tests/test_twin_ladder.py:705-795 (structural tests), 309-333, 638-689, 803-840
% src: C:\source\erisml-compiler-monitor\src\erisml_compiler\runtime\strata.py:105-121
\begin{tikzpicture}[house,
  m/.style={detail, text width=36mm, minimum height=11mm},
  t/.style={detail, text width=50mm, minimum height=11mm},
  c/.style={detail, text width=40mm, minimum height=11mm}]
\node[stepname=Purple] (hm) at (0,0) {Lean model};
\node[stepname=Teal] (ht) at (56mm,0) {What ties it};
\node[stepname=Blue] (hc) at (114mm,0) {Code and scene files};
\foreach \r/\y/\a/\b/\c in {
  1/-9/{\code{permitted}}/{structural tests 1 to 3\\denied by default, guards on system events only, no granting permission}/{\code{margaret\_home.erisml}\\prohibitions and capabilities},
  2/-23/{\code{dispatchPermitted}}/{structural test 4\\the dispatch guard rests on system events only}/{\code{monitoring\_center.erisml}\\dispatch guard},
  3/-37/{\code{Reading.counts}\\the attested flag}/{behaviour tests\\wrong key, altered payload, replay, stale}/{\code{brain.py}\\signature and freshness checks},
  4/-51/{\code{rule}}/{behaviour test\\the governor bar on 0 to 3 witnesses}/{\code{govern\_live}\\in \code{brain.py}},
  5/-65/{\code{authority\_needs\_system}\\\code{absorbing\_stays}}/{loader checks\\a scene that breaks them never runs}/{\code{strata.py}\\\code{Stratification.\_check}}}{
  \node[m] (m\r) at (0,\y mm) {\a};
  \node[t] (t\r) at (56mm,\y mm) {\b};
  \node[c] (c\r) at (114mm,\y mm) {\c};
  \draw[flow, <->] (m\r) -- (t\r);
  \draw[flow, <->] (t\r) -- (c\r);
}
\begin{scope}[on background layer]
  \node[panel=Purple, fit=(hm)(m1)(m5), inner sep=2.5mm] {};
  \node[panel=Teal, fit=(ht)(t1)(t5), inner sep=2.5mm] {};
  \node[panel=Blue, fit=(hc)(c1)(c5), inner sep=2.5mm] {};
\end{scope}
\end{tikzpicture}

  \caption{How each part of the Lean model is tied to the code and the scene files. The scene
  files are checked for the shape the model assumes, and the code is checked by behaviour tests
  and by the loader. None of these ties is itself a proof.}
  \label{fig:lean-bridge}
\end{figure}

The tie is a test, not a proof.
It checks the shape of the scene files and the behaviour of the code on chosen inputs.
It does not prove that the Python interpreter of the scene computes the function $P$.
The model also leaves out one mechanism of the code.
Elevated rights lapse in the code when the fresh witnesses stay below the bar for a set window
(Section~\ref{sec:cont-lapse}).
% src: twin/scene/margaret_home.erisml:805-809
% src: tests/test_twin_ladder.py:1892 (test_rights_revert_when_the_evidence_no_longer_supports_them)
In the model, a restraint ruling once in the history permits restraint for the rest of the run.
% src: formal/twin-containment/TwinContainment.lean:159
A lapse only withdraws permission.
We expect the model to permit at least what the code permits on the same history.
That expectation is not proved.

\section{What the proofs do not cover}\label{sec:lean-notcovered}

The No Escape paper names the attacks a structural containment cannot block.
They are physical attacks on the sensors and political attacks on the people who write the
constraints.
% src: C:\source\erisml-lib\docs\papers\foundations\no_escape.tex:266-274
% src: C:\source\erisml-lib\docs\papers\foundations\no_escape.tex:254-264
The same limits apply here, together with one path the twin leaves open.
% src: formal/twin-containment/TwinContainment.lean:20-22
% src: formal/twin-containment/README.md:31-32

\paragraph{Sensor spoofing below the signature.}
The model trusts the attested flag.
A signature proves that a reading came from a known device and was not altered.
It does not prove that the device measured the world truly.
Someone who holds a lighter under a smoke detector produces a genuine, signed, fresh alert.
Two such sensors on separate substrates produce two witnesses.
The bars raise the cost of such an attack, since it must deceive several independent devices at
once.
They do not remove it.
The No Escape paper classes this as a physical attack, outside its theorem.
% src: C:\source\erisml-lib\docs\papers\foundations\no_escape.tex:274
Trusted sensor hardware is the relevant prior work \cite{liu2012trustedsensors}.

\paragraph{Who writes the scene and the bars.}
The theorems take the scene's permissions and the bars as given.
The model copies its bars from the scene.
The values two and three are choices made by the people who wrote the scene.
The proofs show that the robot respects those values, not that the values are right.
Whether those people and their process are legitimate is a question of governance.
% src: C:\source\erisml-lib\docs\papers\foundations\no_escape.tex:262-264

\paragraph{The centre reads her words with a model.}
One of the three grounds for the centre's dispatch is Margaret's own words, read alone.
% src: formal/twin-containment/TwinContainment.lean:408-414
A model reads them.
A misreading of her words as a request for help could therefore send emergency services.
% src: formal/twin-containment/README.md:31-32
The robot's message cannot reach this path, as \code{dispatch\_ignores\_the\_robot} shows.
A false reading of her own speech can.
This is the one remaining path by which a model output reaches authority in the twin.
The consequence is a dispatch to someone who may not need it, not force against her.

\paragraph{The model is not the code.}
The structural tests check the scene's shape.
They do not verify the interpreter, the cryptography, the operating system or the robot's
actuators.
Lapse of rights is outside the model, as noted above.
The proofs also say nothing about whether the robot's permitted choices are good ones.
They bound what the robot may do.
They do not rank what it should do.
Dennis and colleagues proved properties of ethical choice in an agent's plan selection
\cite{dennis2016formal}.
The Lean results here are of a different kind.
They concern which events can move authority, not which permitted plan is best.
